% !PRINTED 249
\weberpage{249}

Endlich ist noch die Bedingung für die Gleichheit aller vier Wurzeln
\weq{13}{a = 0, \quad b = 0, \quad c = 0.}

Man kann diese Verhältnisse sehr anschaulich machen durch eine geometrische Deutung, und wenn auch geometrische Betrachtungen nicht eigentlich in unserem Plane liegen, so wollen wir doch nicht unterlassen, den Leser darauf gelegentlich hinzuweisen.

Deutet man $a$, $b$, $c$ als rechtwinklige Coordinaten im Raume, so ist jeder Raumpunkt als Träger einer gewissen biquadratischen Gleichung von der Form (1), $f(x) = 0$, zu betrachten; alle Gleichungen, die eine bestimmte Zahl $x$ zur Wurzel haben, werden durch Punkte einer Ebene $[f(x) = 0]$ repräsentirt. Die Gleichung $D = 0$ ist die Gleichung einer krummen Oberfläche (fünften Grades), der Discriminantenfläche, die von den Schnittlinien der Ebenen $f(x) = 0$, $f'(x) = 0$ erzeugt wird, und also eine abwickelbare Fläche ist. Sie ist die Einhüllende aller Ebenen $f(x) = 0$.

Die Fläche hat eine aus zwei Zweigen bestehende Rückkehrkante, die durch die Gleichungen (11), (12) dargestellt ist, und eine Doppellinie, die durch die Gleichungen $b = 0$, $a^2 - 4c = 0$ bestimmt ist und also die Gestalt einer Parabel hat. Auf der Seite der negativen $a$ durchsetzt sich in dieser Parabel die Fläche selbst. Auf der Seite der positiven $a$ setzt sich die Parabel als isolirte Linie fort.

Die Discriminantenfläche theilt den ganzen Raum in drei Fächer, von denen zwei nur längs der Doppelparabel zusammenhängen, und diese Fächer enthalten die Punkte, denen keine, zwei und vier reelle Wurzeln entsprechen. Nennen wir für den Augenblick diese drei Fächer $[0]$, $[2]$, $[4]$, so grenzt $[0]$ an $[4]$ nur längs der Doppelparabel, während $[0]$ sowohl als $[4]$ längs der Flächentheile an $[2]$ grenzen.

% [sic] The source defines [4,2] but here prints the second boundary pair as [0],[4]; preserved diplomatically.
Wir wollen mit $[0, 2]$ und $[4, 2]$ die Grenzflächen von $[0]$, $[2]$ und von $[0]$, $[4]$ bezeichnen. Der isolirte Theil der Parabel setzt sich in das Innere von $[0]$ fort. Die Rückkehrkanten liegen auf dem Theil der Fläche, der $[4]$ von $[2]$ scheidet.

Im Coordinatenanfang stossen die drei Raumtheile und ihre Grenzflächen und Grenzlinien zusammen. Die Punkte der Flächentheile $[0, 2]$ stellen Gleichungen mit zwei gleichen und
%CONT>>
